\section{Introduction}
Grid-based kinetic plasma simulations based on the Vlasov equation evolve the distribution function directly in phase space rather than sampling it with numerical particles. In contrast to particle-in-cell (PIC) methods, the continuum approach is free of sampling noise and therefore gives direct access to fine distribution-function structure, weak phase-space features, and wall fluxes. For that reason, continuum kinetic methods are especially attractive in bounded plasma problems where sheath structure, non-Maxwellian features, and plasma-material interaction must be resolved with high fidelity. 

One of the main obstacles for applying the Vlasov method in the real engineering is geometric complexity. In a Vlasov solver, the physical-space mesh is coupled to a discretization in velocity space, and every additional cost in configuration space is effectively replicated across the velocity grid. As a result, design choices that are common in CFD, such as body-fitted meshes, heavy local refinement near curved walls, or distinct meshes for transport and field equations, can become much more expensive and much harder to manage in the grid-based kinetic computations. Recent high-order GPU work on Vlasov-Poisson solvers again highlights that the very large number of phase-space unknowns is a central practical limitation of the grid-based kinetic methods.

One possible approach for resolving this issue is to use body-fitted or unstructured meshes directly in phase space. This direction was explored early, for example, by Besse and Sonnendrücker, who developed semi-Lagrangian Vlasov schemes on unstructured phase-space meshes. However, the literature also makes clear why this route is difficult to scale in arbitrary geometry. Filbet and Yang explicitly argued that, because of the high-dimensional nature of kinetic equations, unstructured-mesh treatments are often unattractive and that Cartesian-grid techniques are preferable when possible. For Vlasov solvers, this argument is particularly compelling because many high-order Eulerian and semi-Lagrangian algorithms derive much of their efficiency from structured-grid operations.

On the other hand, a uniform mesh with arbitrary boundary condition can be a promising approach. More recently, Liu et al. developed an immersed finite element and conservative semi-Lagrangian framework for plasma-material interaction on structured interface-independent meshes, combining a first order half-way ghost-cell treatment for Vlasov transport with an immersed field solve. These works clearly demonstrate progress, but they also show that the design space remains open, particularly for time-dependent, same-mesh, sharp-interface Vlasov–Poisson formulations with explicit near-boundary accuracy analysis.

A closely related body of work exists in deterministic rarefied-gas simulation, where discrete-velocity and BGK/Boltzmann solvers face a very similar tension between velocity-space cost and geometric fidelity. There, Cartesian embedded-boundary ideas have proved especially attractive. Mieussens’ survey emphasized immersed Cartesian techniques as a promising direction for deterministic solvers. Filbet and Yang developed an inverse Lax-Wendroff boundary treatment for Boltzmann-type models in arbitrary geometry on Cartesian meshes. Dechriste and Mieussens later showed that conservative Cartesian cut-cell methods can combine geometric flexibility, exact conservation, and a relatively simple implementation for rarefied flows around moving obstacles. This rarefied-gas literature strongly suggests that a Cartesian sharp-interface treatment is a natural strategy for grid-based Vlasov solvers as well.

The contribution of the present work is therefore not simply to introduce complex geometry into Vlasov simulation, since that topic has already seen meaningful progress, but to provide a unified same-mesh formulation in which both the Vlasov and Poisson equations are treated with second-order sharp-interface accuracy near the boundary. This is particularly useful when the interior Vlasov solver is of higher order: a second-order near-boundary closure is accurate enough to remain consistent with the interior discretization on practical meshes, while avoiding the very dense local refinement that would otherwise be required around curved objects. 

The remainder of the paper is organized as follows. Section 2 introduces the mathematical model and boundary conditions. Section 3 presents the unified same-mesh numerical method. Section 4 discusses the consistency and accuracy analysis near the immersed boundary. Section 5 contains numerical verification and validation. Section 6 presents representative application cases. Section 7 concludes the paper.
